% GENERATED FILE. Edit canonical lessons, then run npm run generate:books.
% canonical-source-hash: fnv1a64:bf7a3bb06fe8bbb8
% canonical-lessons: GU-C149-bolish, GU-C149-bolshe, GU-C149-jaish, GU-C149-ungh, GU-W48-uu, GU-C149-future-practice

\chapter{What will happen}
\label{ch:gu-what-will-happen}
\hlchaptermodality{149}

\begin{chapteropening}
\textbf{By the end of this chapter:} \emph{I can say what I and other people will do tomorrow with the future endings, ask when someone will come, and say that I am sleepy.}

Tell tomorrow from morning to night, then ask someone when they will come and say where you will be.
\end{chapteropening}

\section[\texorpdfstring{…\gu{ીશ} · …\gu{ીશું} (-īś · -īśũ)}{…ીશ · …ીશું (-īś · -īśũ)}]{\gu{હું} \gu{બોલીશ} — the future, built on \gu{મળીશું}}

\label{lesson:GU-C149-bolish}



\begin{quote}
Say \emph{I drank tea} and \emph{I ate rice}. Then the goodbye from the first week: \textbf{\gu{કાલે} \gu{મળીશું}}, \emph{see you tomorrow}, word for word \emph{tomorrow we will meet}.
\end{quote}

\begin{grammarlens}[title={Grammar Lens: -\gu{ીશું}, we will}]
\textbf{\gu{મળીશું}} is \textbf{\gu{મળ}-}, \emph{meet}, plus \textbf{-\gu{ીશું}}, \emph{we will}. Gujarati needs no separate word for \emph{will}: the future is an ending on the stem.

\begin{quote}
\textbf{\gu{આપણે} \gu{કાલે} \gu{મળીશું}.} — \emph{āpṇe kāle maḷīśũ.} — We'll meet tomorrow. \textbf{\gu{અમે} \gu{ગુજરાતી} \gu{બોલીશું}.} — \emph{ame gujarātī bolīśũ.} — We will speak Gujarati.
\end{quote}
\end{grammarlens}

\subsection*{You'll want to know: -\gu{ીશ}, I will}
Take the last \textbf{-\gu{ું}} off and you have \emph{I will}:

\begin{quote}
\textbf{\gu{હું} \gu{ગુજરાતી} \gu{બોલીશ}.} — \emph{hũ gujarātī bolīś.} — I will speak Gujarati.
\end{quote}

The familiar \emph{you} takes the same form in the standard language: \textbf{\gu{તું} \gu{બોલીશ}}, \emph{you will speak}. There is no copula and no gender: a man and a woman both say \textbf{\gu{હું} \gu{બોલીશ}}.

The \textbf{-\gu{શ}} of these endings is usually traced to the old Sanskrit future, the \emph{-iṣya-} of \emph{kariṣyati}, \textquotedblleft{}will do\textquotedblright{}.

\subsection*{Your turn}
Say these aloud:

\begin{itemize}
\raggedright
  \item \emph{hũ bolīś}, \emph{ame bolīśũ}
  \item \textbf{\gu{હું} \gu{આવતીકાલે} \gu{લખીશ}} — \emph{hũ āvtīkāle lakhīś}, I'll write tomorrow
  \item \textbf{\gu{અમે} \gu{ઘરે} \gu{રહીશું}} — \emph{ame ghare rahīśũ}, we'll stay at home
  \item \textbf{\gu{અમે} \gu{કાલે} \gu{આવીશું}} — \emph{ame kāle āvīśũ}, we'll come tomorrow
\end{itemize}

\begin{tcolorbox}[breakable,colback=teal!4,colframe=teal!35!black,title={Before you move on}]
What word for \emph{will} does Gujarati use? (None; the future is an ending.) Which ending did \textbf{\gu{મળીશું}} teach you long ago? (\textbf{-\gu{ીશું}}, \emph{we will}.)
\end{tcolorbox}
\section[\texorpdfstring{…\gu{શે} · …\gu{શો} (-śe · -śo)}{…શે · …શો (-śe · -śo)}]{\gu{તે} \gu{બોલશે} — the future for everyone else}

\label{lesson:GU-C149-bolshe}



\begin{quote}
Say \emph{I will speak Gujarati}, then \emph{we'll meet tomorrow}.
\end{quote}

\begin{grammarlens}[title={Grammar Lens: -\gu{શે}}]
\begin{quote}
\textbf{\gu{તે} \gu{આવતીકાલે} \gu{આવશે}.} — \emph{te āvtīkāle āvśe.} — He will come tomorrow. She will come tomorrow. \textbf{\gu{તેઓ} \gu{ગુજરાતી} \gu{બોલશે}.} — \emph{teo gujarātī bolśe.} — They will speak Gujarati.
\end{quote}

\textbf{-\gu{શે}} sits on the bare stem, with no \textbf{-\gu{ી}-} before it: \textbf{\gu{બોલ}-\gu{શે}}, \textbf{\gu{આવ}-\gu{શે}}.
\end{grammarlens}

\subsection*{You'll want to know: -\gu{શો}}
\textbf{\gu{તમે}} keeps the \textbf{-\gu{ો}} it has in the present, \textbf{\gu{બોલો} \gu{છો}}:

\begin{quote}
\textbf{\gu{તમે} \gu{ક્યારે} \gu{આવશો}?} — \emph{tame kyāre āvśo?} — When will you come?
\end{quote}

The future for one person:

\noindent
\begin{tabularx}{\linewidth}{@{}>{\raggedright\arraybackslash}X>{\raggedright\arraybackslash}X@{}}
\toprule
\textbf{who} & \textbf{will speak} \\
\midrule
\gu{હું}, \gu{તું} & \gu{બોલીશ} \\
\gu{તે} & \gu{બોલશે} \\
\bottomrule
\end{tabularx}

And for more than one, or for respect:

\noindent
\begin{tabularx}{\linewidth}{@{}>{\raggedright\arraybackslash}X>{\raggedright\arraybackslash}X@{}}
\toprule
\textbf{who} & \textbf{will speak} \\
\midrule
\gu{અમે}, \gu{આપણે} & \gu{બોલીશું} \\
\gu{તમે} & \gu{બોલશો} \\
\gu{તેઓ} & \gu{બોલશે} \\
\bottomrule
\end{tabularx}

\subsection*{Your turn}
\begin{itemize}
\raggedright
  \item \emph{Say it:} \emph{te āvśe}, then \emph{teo āvśe}
  \item \emph{Ask:} \emph{tame kyāre āvśo?}
  \item \emph{Answer:} with \textbf{\gu{હું} … \gu{આવીશ}} — \emph{hũ … āvīś}
\end{itemize}

\begin{tcolorbox}[breakable,colback=teal!4,colframe=teal!35!black,title={Before you move on}]
Which two persons share \textbf{-\gu{શે}}? (\textbf{\gu{તે}} and \textbf{\gu{તેઓ}}.) What does \textbf{\gu{તમે}} take? (\textbf{-\gu{શો}}.)
\end{tcolorbox}
\section[\texorpdfstring{\gu{જઈશ} · \gu{લઈશ} · \gu{થશે} … (jaīś · laīś …)}{જઈશ · લઈશ · થશે … (jaīś · laīś …)}]{\gu{જઈશ}, \gu{થશે}, \gu{હશે} — the short verbs}

\label{lesson:GU-C149-jaish}



\begin{quote}
Say \emph{she will come tomorrow}, then ask \emph{when will you come?} And two short verbs: \textbf{\gu{જવું}}, \emph{to go}, and \textbf{\gu{લેવું}}, \emph{to take}.
\end{quote}

\begin{grammarlens}[title={Grammar Lens: \gu{જવું}, \gu{લેવું}, \gu{કરવું}}]
\textbf{\gu{જવું}} and \textbf{\gu{લેવું}} have short stems, and they put an \textbf{\gu{ઈ}} where the long verbs have \textbf{-\gu{ી}-}:

\begin{itemize}
\raggedright
  \item \textbf{\gu{હું} \gu{જઈશ}} — \emph{jaīś} — I will go; \textbf{\gu{અમે} \gu{જઈશું}}; \textbf{\gu{તે} \gu{જશે}}; \textbf{\gu{તમે} \gu{જશો}};
  \item \textbf{\gu{હું} \gu{લઈશ}} — \emph{laīś} — I will take; \textbf{\gu{તે} \gu{લેશે}}.
\end{itemize}

\textbf{\gu{કરવું}} is regular: \textbf{\gu{હું} \gu{કરીશ}}, \textbf{\gu{અમે} \gu{કરીશું}}, \textbf{\gu{તે} \gu{કરશે}}.

\begin{quote}
\textbf{\gu{હું} \gu{આવતીકાલે} \gu{બજાર} \gu{જઈશ}.} — \emph{hũ āvtīkāle bajār jaīś.} — I'll go to the market tomorrow.
\end{quote}
\end{grammarlens}

\subsection*{You'll want to know: \gu{થશે} and \gu{હશે}}
\textbf{\gu{થવું}}, \emph{to happen}, and \textbf{\gu{હોવું}}, \emph{to be}, are heard most often with \textbf{-\gu{શે}}:

\begin{quote}
\textbf{\gu{મોડું} \gu{થશે}.} — \emph{moḍũ thaśe.} — It will be late, word for word \emph{lateness will happen}. \textbf{\gu{તે} \gu{ઘરે} \gu{હશે}.} — \emph{te ghare haśe.} — She will be at home.
\end{quote}

\textbf{\gu{હશે}} is the future of \textbf{\gu{છે}}: \textbf{\gu{છે}} now, \textbf{\gu{હતો}} before, \textbf{\gu{હશે}} later. With \textbf{\gu{તમે}} it is \textbf{\gu{હશો}}: \textbf{\gu{શું} \gu{તમે} \gu{ઘરે} \gu{હશો}?}, \emph{will you be at home?} Said about right now, \textbf{\gu{હશે}} is a guess: \textbf{\gu{તે} \gu{ઘરે} \gu{હશે}} can also mean \emph{she must be at home}.

\subsection*{Your turn}
Say these aloud:

\begin{itemize}
\raggedright
  \item \emph{hũ jaīś}, \emph{te jaśe}, \emph{tame jaśo}
  \item \textbf{\gu{હું} \gu{દવા} \gu{લઈશ}} — \emph{hũ davā laīś}
  \item \textbf{\gu{સાંજે} \gu{તેઓ} \gu{ઘરે} \gu{હશે}} — \emph{sānje teo ghare haśe}
\end{itemize}

\begin{tcolorbox}[breakable,colback=teal!4,colframe=teal!35!black,title={Before you move on}]
Which form is the future of \textbf{\gu{છે}}? (\textbf{\gu{હશે}}.) What do \textbf{\gu{જવું}} and \textbf{\gu{લેવું}} put before \textbf{-\gu{શ}}? (\textbf{\gu{ઈ}}: \textbf{\gu{જઈશ}}, \textbf{\gu{લઈશ}}.)
\end{tcolorbox}
\section[\texorpdfstring{\gu{ઊંઘ} (ūṅgh)}{ઊંઘ (ūṅgh)}]{\gu{ઊંઘ} (ūṅgh) — sleepiness, sleep}

\label{lesson:GU-C149-ungh}



\begin{quote}
Before the new one: say \emph{I will go}, then \emph{it will be late}.
\end{quote}

\subsection*{You'll want to know: \gu{ઊંઘ}}
\textbf{\gu{ઊંઘ}} (\emph{ūṅgh}) — \textquotedblleft{}sleepiness, sleep\textquotedblright{}.

Gujarati does not say \emph{I am sleepy}. Sleep comes \emph{to} you, the way a fever does:

\begin{quote}
\textbf{\gu{મને} \gu{ઊંઘ} \gu{આવે} \gu{છે}.} — \emph{mane ūṅgh āve che.} — I'm sleepy. Word for word: \emph{to me, sleep comes}. \textbf{\gu{આજે} \gu{રાત્રે} \gu{સારી} \gu{ઊંઘ} \gu{આવશે}.} — \emph{āje rātre sārī ūṅgh āvśe.} — Tonight there will be good sleep.
\end{quote}

\textbf{\gu{ઊંઘ}} is feminine, so \emph{good} bends to \textbf{\gu{સારી}}. \textbf{\gu{સૂવું}}, \emph{to sleep}, is lying down to sleep; \textbf{\gu{ઊંઘ}} is the sleep itself.

\begin{cousinweb}
\textbf{\gu{ઊંઘ}} (\emph{ūṅgh}) is inherited. CDIAL reconstructs an Indo-Aryan \emph{uṅghā}, \textquotedblleft{}sleep\textquotedblright{}, and lists the Gujarati word under it, beside Hindi \emph{ūṅgh} (CDIAL 1633).
\end{cousinweb}

\begin{grammarlens}[title={What you've built}]
One more word, and sleep joins likes and fevers among the things that happen \emph{to} someone.
\end{grammarlens}

\subsection*{Your turn}
Say these aloud:

\begin{itemize}
\raggedright
  \item \emph{ūṅgh}
  \item \emph{mane ūṅgh āve che}
  \item \textbf{\gu{શું} \gu{તમને} \gu{ઊંઘ} \gu{આવે} \gu{છે}?} — \emph{śũ tamne ūṅgh āve che?}, are you sleepy?
\end{itemize}

\begin{tcolorbox}[breakable,colback=teal!4,colframe=teal!35!black,title={Before you move on}]
What does \textbf{\gu{ઊંઘ}} mean? (\emph{Sleepiness, sleep}.) Which letter in it have you read but never written? (\textbf{\gu{ઊ}}, the next lesson.)
\end{tcolorbox}
\section[\texorpdfstring{\gu{ઊ} (ū)}{ઊ (ū)}]{\gu{ઊ} — the standing vowel ū}

\label{lesson:GU-W48-uu}



\begin{quote}
\begin{itemize}
\raggedright
  \item \emph{Recall:} say \emph{sleepiness} — \textbf{R1}, one lesson back
  \item \emph{Recall:} say the standing vowel \textbf{\gu{ઉ}}, and the sign \textbf{\gu{ૂ}} that gives a consonant a long \emph{ū}
\end{itemize}

\textbf{\gu{ઊંઘ}} opens on a long \emph{ū} with no consonant to carry it.
\end{quote}

\subsection*{You'll want to know: \gu{ઉ} with a long tail}
\noindent
\begin{tabularx}{\linewidth}{@{}>{\raggedright\arraybackslash}X>{\raggedright\arraybackslash}X>{\raggedright\arraybackslash}X@{}}
\toprule
\textbf{} & \textbf{after a consonant} & \textbf{standing alone} \\
\midrule
short \emph{u} & \textbf{\gu{ુ}} & \textbf{\gu{ઉ}} \\
long \emph{ū} & \textbf{\gu{ૂ}} & ? \\
\bottomrule
\end{tabularx}

The missing letter is built on \textbf{\gu{ઉ}}: the whole of \textbf{\gu{ઉ}}, then a tail that runs across the top and down the right side.

\subsection*{Script — meet one form}
\begin{quote}
\textbf{\gu{ઊ}} — \emph{ū}
\end{quote}

Keep \textbf{\gu{ઊ}} visible. Write \textbf{\gu{ઉ}} first: the small upper bowl, the middle cusp, the broad lower bowl, and the tall curve on the outside. Do not lift. Carry on across the upper shoulder, come down the long tail on the right, and finish through its foot. One stroke, no lift.

The tail is the whole difference. Stop before it and the letter reads \textbf{\gu{ઉ}}.

\subsection*{Writing — notice and trace}
\hlblockfigure{\detokenize{figures/GU-W48-uu-filmstrip.pdf}}{How \gu{ઊ} is written, stroke by stroke}

Finger-trace \textbf{\gu{ઊ}} once, then air-trace it once. Notice where \textbf{\gu{ઉ}} ends and the tail begins.

\subsection*{Writing — copy with the model}
Copy visible \textbf{\gu{ઊ}} once. Compare only two things: is all of \textbf{\gu{ઉ}} there, and does the tail reach down the right side? Repair one landmark if needed.

\subsection*{Writing — recall after hiding}
Hide the model and write the standing vowel \emph{ū} once. Uncover \textbf{\gu{ઊ}}, compare, and stop after one repair.

\begin{tcolorbox}[breakable,colback=teal!4,colframe=teal!35!black,title={Before you move on}]
Read \textbf{\gu{ઉ} \gu{ઊ} \gu{ઊંઘ}} aloud. What turns \textbf{\gu{ઉ}} into \textbf{\gu{ઊ}}? (\textbf{The long tail.})

Sources: \href{https://www.t30apps.com/gujarati-alphabet-writing/}{Gujarati Alphabet Writing Practice} and the \href{https://www.unicode.org/charts/PDF/U0A80.pdf}{Unicode Gujarati chart}.
\end{tcolorbox}
\section[Practice]{Tomorrow — a practice}

\label{lesson:GU-C149-future-practice}



\begin{quote}
Say \emph{I'm sleepy}.

\begin{itemize}
\raggedright
  \item \emph{Write it:} \textbf{\gu{ઊ}} once from memory
\end{itemize}
\end{quote}

\begin{tcolorbox}[breakable,skin=enhanced,colback=orange!6,colframe=orange!45!black,boxrule=0.5pt,arc=1mm,left=6pt,right=6pt,top=4pt,bottom=4pt,fonttitle=\bfseries,title={Read this}]
\begin{quote}
\textbf{\gu{આવતીકાલે} \gu{સવારે} \gu{હું} \gu{વહેલો} \gu{ઉઠીશ}.} \textbf{\gu{બપોરે} \gu{હું} \gu{બજાર} \gu{જઈશ}.} \textbf{\gu{સાંજે} \gu{મારી} \gu{બહેન} \gu{આવશે}.} \textbf{\gu{અમે} \gu{ઘરે} \gu{ચા} \gu{પીશું}.} \textbf{\gu{રાત્રે} \gu{મને} \gu{ઊંઘ} \gu{આવશે}.}
\end{quote}

Say these aloud:

\begin{itemize}
\raggedright
  \item the five lines aloud
  \item the same plan for your sister: \emph{te … uṭhśe}, \emph{te … jaśe}
\end{itemize}
\end{tcolorbox}

\subsection*{Your turn}
\begin{itemize}
\raggedright
  \item \emph{Answer:} \textbf{\gu{તમે} \gu{આવતીકાલે} \gu{શું} \gu{કરશો}?} — \emph{tame āvtīkāle śũ karśo?}
  \item \emph{Say it:} \textbf{\gu{મોડું} \gu{થશે}} — \emph{moḍũ thaśe} — and say whose tomorrow it might spoil
\end{itemize}

\begin{tcolorbox}[breakable,colback=teal!4,colframe=teal!35!black,title={Before you move on}]
Which line told someone else's future? (\textbf{\gu{સાંજે} \gu{મારી} \gu{બહેન} \gu{આવશે}}.) Which verb put an \textbf{\gu{ઈ}} before its ending? (\textbf{\gu{જવું}}: \textbf{\gu{જઈશ}}.)
\end{tcolorbox}
